\begin{table}[H]
\centering
\caption{Plataformas de validación empírica declaradas en el corpus (n=80). Un estudio puede declarar más de una plataforma.}
\label{tab:rq3_plataformas}
\small
\renewcommand{\arraystretch}{1.3}
\begin{tabular}{|>{\raggedright\arraybackslash}m{0.44\linewidth}|c|c|}
\hline
\textbf{Plataforma de validación} & \textbf{Estudios} & \textbf{\% del corpus} \\
\hline
FPGA & 8 & 10.0\% \\
\hline
NVIDIA Jetson & 4 & 5.0\% \\
\hline
Raspberry Pi / SBC & 2 & 2.5\% \\
\hline
NPU / ASIC / acelerador dedicado & 1 & 1.2\% \\
\hline
SoC móvil comercial (Snapdragon, Apple M1) & 4 & 5.0\% \\
\hline
CPU ARM / RISC-V & 1 & 1.2\% \\
\hline
GPU escritorio/servidor \textit{(fuera de alcance Edge)} & 6 & 7.5\% \\
\hline
\textbf{Sin plataforma física identificable} & \textbf{57} & \textbf{71.2\%} \\
\hline
\multicolumn{3}{|l|}{\footnotesize Estudios con validación verificable en hardware \emph{Edge} concreto: \textbf{18/80 (22\%)}} \\
\hline
\end{tabular}
\end{table}